\documentclass[12pt]{article}

\usepackage[a4paper, total={6in, 8in}]{geometry}
\usepackage{textcomp}

\begin{document}
\setlength{\parindent}{0pt}

\def \t {\textrightarrow}
\def \v {\vspace{1em}}
\def \bi {\begin{itemize}}
\def \ei {\end{itemize}}
\def \s[#1] {\section*{#1}}
\def \ss[#1] {\subsection*{#1}}
\def \sss[#1] {\subsubsection*{#1}}

\s[Riprendere appunti vecchi (se ce ne sono forse)]
\s[Quintilliano]
Si parla di impero, ruolo del docente \t il suo ruolo viene ricordato come quello di un docente pubblico \t di scuola pubblica \\
Il docente acquisisce una sua fisionomia autonoma e viene stipendiato \t e nasce il concetto di classe (come gruppo di studenti) \\
L'istruzione domiciliare viene sottoposta a un confronto con l'istruzione uno a molti (in classe) \t dinamiche comunicative diverse \t tutto questo accompagnato da uno studio in merito all eta dello studente \\
Lo studente nel mondo romano era considerato in adolescenza fino ai 30 anni \t il cursus onorum si chiude a questa eta \\
Questo si lega con la con il fanciullino di Pascoli

\v

"Istitutio oratoria" \t chiama in causa il rapporto tra l'intellettuale e il potere \t che si misurano sulla base di un equilibrio che fa guardagnare terreno all'intellettuale solo se la sua formazione politica avviene fino da piccolo \\
Quindi la preparazione politica è importante, dove viene fuori il confronto con Cicerone \t cicerone era il predecessore, e aveva parlato di una preparazione dell'oratore che doveva essere \textbf{vir bonus dicendi periturs} METTERLO NEL CONPITO \\
L'oratore deve docere, movere, flectere (cicerone) \\
Lo studente sara il futuro dirigente della roma \t e deve capire se vale di piu l istruzione domiciliare con il precettore o quella scolastica con il maestro? \t il maestro deve essere emotivamente coinvolto, in modo da sviluppare empatia e stimolare lo studente \\
La paura da socialita \t fobia del sociale

\v

Inoltre il precettore se lo potevano permettere solo persone + abbienti \t quindi selettivita dell'istruzione \t ma inoltre questa istruzione 1 a 1, è deleteria per la crescita \t il bambino non entra in contatto con altri bambini \\
Nella socialita il bambino è facilitato a uscire dalla fase anale (freud) \\
La scula intesa come gruppo classe era prevista per una famiglia in particolare mononucleare, senza fratelli \\
Questo esperimento si manifesta quando ci sono famiglie nucelarei o quando il livello di sdislivello tra i fratelli è mollto alto

\v

Cicerone afferma che l oratore è un fact totum, un eclettico \t ma qua in quintilliano viene mantenuto l'interesse per la concinnitas (equilibirio nel discorso, sic ut ita ut, tam ut, sic sic = relazione tra le parti) \\
Quintilliano pero non ha la stessa vis dell oratoria di cicerone, perche sotto l impero non puo esprimersi come vuole \t la maglia ingabbiante del potere è presente \\
Questo l'abbiamo visto in Fedro, e Seneca
\end{document}
